% ==================== DISEÑO DE LA SOLUCIÓN ====================
\section{Diseño de la Solución}

\subsection{Escenario de Aplicación}

El escenario de trabajo lo integra un equipo \host{} que ejecuta GNU/Linux con Oracle VM VirtualBox~7.2 y con la solución en Bash instalada, junto con dos máquinas virtuales GNU/Linux llamadas \texttt{VM1} y \texttt{VM2}, que comparten el mismo sistema invitado, Debian~13 (trixie), pero se configuran con parámetros distintos. El \host{} participa en la red local a través de una de sus interfaces físicas. En este escenario esa interfaz es \texttt{enp2s0}, aunque el nombre depende de la topología de cada equipo. Cada \vm{} posee una interfaz de red en modo puente que la conecta a la red local \texttt{192.168.1.0/24} y que en el invitado se denota \texttt{enp0s3}. Tal designación resulta del controlador y del orden de los adaptadores, de modo que, igual que las direcciones, constituye un valor del escenario y no un dato del que la solución dependa. En el modo puente la dirección tampoco la fija la \vm{}, sino que la ofrece el servidor DHCP de la red local, de modo que no puede predecirse a partir de la configuración del escenario \parencite{virtualbox_networking}. El entorno de pruebas asigna \texttt{192.168.1.51} a \texttt{VM1}, \texttt{192.168.1.52} a \texttt{VM2} y \texttt{192.168.1.10} al \host{}. Ese plan describe únicamente el entorno y no constituye una premisa de la solución, la cual identifica la dirección en tiempo de ejecución, tal como se detalla en la Subsección~\ref{sub:arquitectura}. Asimismo, ambas cuentan con un disco virtual adicional (\texttt{.vdi}) que la automatización deja montado como almacenamiento de datos: \texttt{/mnt/datos} con ext4 en la primera y \texttt{/srv/datos} con xfs en la segunda \parencite{linux_storage_management}. La comunicación entre el \host{} y cada \vm{} se realiza con la Guest Control API, que se apoya en las \guestadditions{} instaladas en el invitado \parencite{virtualbox_guestcontrol}.

La Figura~\ref{fig:escenario} representa el escenario mediante un diagrama de despliegue UML elaborado con PlantUML, en el que se distinguen la interfaz física del \host{}, la red local en modo puente, la interfaz de cada \vm{}, los discos de almacenamiento y el flujo general de la automatización en cinco pasos: (1) la solución enciende la \vm{} con \texttt{VBoxManage startvm}; (2) identifica en tiempo de ejecución la dirección IP de la \vm{} a partir de la propiedad que publican las \guestadditions{}, con respaldo sobre la tabla ARP de la red en puente; (3) ejecuta comandos en el invitado con \texttt{VBoxManage guestcontrol}; (4) prepara el almacenamiento mediante la secuencia de identificación del disco, particionado, creación del sistema de archivos, punto de montaje, entrada en \texttt{/etc/fstab} y montaje; y (5) verifica el resultado con \texttt{findmnt} y \texttt{df -h}, incluida la re-ejecución idempotente sobre una \vm{} ya configurada \parencite{virtualbox_guestadditions, linux_storage_management, idempotency_concept}. El mismo flujo se aplica sobre \texttt{VM2} con otros parámetros, lo que evidencia la reutilización de la solución. Las direcciones IP y los nombres de interfaz que aparecen en la figura describen el entorno de pruebas y no son datos de entrada de la solución: la dirección la obtiene ésta por sí misma en el paso~2 y de los nombres de interfaz no hace uso alguno.

\begin{figure}[H]
    \centering
    \includegraphics[width=0.95\textwidth]{escenario.png}
    \caption{Escenario de despliegue: \host{}, VirtualBox, red en modo puente, \vm{} y almacenamiento adicional}
    \label{fig:escenario}
\end{figure}

\subsection{Arquitectura y Comunicación entre Host y Guest}
\label{sub:arquitectura}

La solución se estructura en tres planos diferenciados. Un punto de entrada único, \texttt{main.sh}, recibe el nombre de la \vm{} objetivo, carga su archivo de configuración desde \texttt{config/} y coordina la ejecución; las bibliotecas de \texttt{lib/} concentran funciones reutilizables de orquestación, comunicación y almacenamiento; y los archivos de configuración por \vm{} declaran los parámetros que cambian entre máquinas: disco objetivo, punto de montaje, sistema de archivos y credenciales \parencite{bash_best_practices}. Con esta separación, el mismo algoritmo se aplica a \texttt{VM1} y \texttt{VM2} sin modificar el código, y las verificaciones con ShellCheck y BATS descritas en el Marco Teórico recaen sobre funciones pequeñas e independientes \parencite{shellcheck, bats_testing}. El lenguaje de implementación es Bash, justificado en la Subsección~\ref{sub:bash} por su disponibilidad nativa en el \host{} y por su acceso directo a la interfaz de línea de comandos \texttt{VBoxManage} \parencite{virtualbox_vboxmanage}.

La comunicación con el invitado se establece con la Guest Control API (\texttt{VBoxManage guestcontrol}), mecanismo que satisface el requerimiento de utilizar las \guestadditions{} para la interacción \parencite{virtualbox_guestcontrol}. En este mecanismo el \host{} entrega las órdenes al hipervisor, éste las transporta al servicio \texttt{VBoxService} en ejecución dentro de la \vm{} y el resultado regresa por el mismo canal; ese trayecto no atraviesa la red en puente, de modo que la ejecución de comandos no depende de conocer de antemano la dirección IP del invitado \parencite{virtualbox_guestadditions, virtualbox_advanced_topics}. La autenticación se realiza con un usuario y una contraseña del propio sistema invitado, leídos del archivo de configuración antes de abrir la sesión, de manera que ninguna etapa solicita credenciales, comandos ni pulsaciones al usuario \parencite{virtualbox_guestcontrol}.

La dirección IP de la \vm{} sí se identifica de forma automática, porque es un requisito del proceso y se emplea para verificar la conectividad y registrar la máquina atendida. El mecanismo principal consulta la propiedad que publican las \guestadditions{} en el hipervisor. La orden \texttt{VBoxManage guestproperty get VM1} lee la propiedad \path{/VirtualBox/GuestInfo/Net/0/V4/IP} y devuelve la dirección vigente, sin sondear la red ni exigir servicios adicionales en el invitado \parencite{virtualbox_guestadditions, virtualbox_vboxmanage}. Como respaldo, y en particular si la propiedad aún no se ha actualizado tras el arranque, la solución consulta la tabla ARP completa del \host{} con \texttt{ip neigh} y selecciona la entrada cuya dirección MAC coincide con la del adaptador de la \vm{} reportada por \texttt{VBoxManage showvminfo}, operación que no exige conocer el nombre de ninguna interfaz; si éste fuera necesario, se obtiene en tiempo de ejecución con \texttt{ip route get} a partir de la dirección recién identificada. Del mismo modo, la solución nunca utiliza el nombre de la interfaz del invitado, el cual ni la Guest Control ni la consulta de propiedades requieren. En ambos casos la dirección IP es un resultado que la solución obtiene, valida y registra, nunca un valor fijo codificado en el programa \parencite{virtualbox_networking}.

\subsection{Algoritmo de Automatización (Idempotente)}
\label{sub:algoritmo}

El algoritmo se organiza en tres etapas que se ejecutan de forma secuencial y sin intervención del usuario: una etapa de preparación en el \host{}, una etapa de preparación del almacenamiento que se despacha al invitado con \texttt{VBoxManage guestcontrol} y una etapa de verificación y cierre que vuelve a ejecutarse en el \host{}. La Figura~\ref{fig:algoritmo} resume ese flujo en un diagrama de procesos con dos carriles, lo que permite distinguir en cada nodo qué equipo ejecuta la acción.

En la etapa de preparación, el punto de entrada carga la configuración de la \vm{} objetivo, comprueba si está encendida y, de no estarlo, la enciende con \texttt{VBoxManage startvm} y aguarda la respuesta de \texttt{VBoxService}; después identifica su dirección IP en tiempo de ejecución, según se describió en la Subsección~\ref{sub:arquitectura}, y despacha las órdenes al invitado \parencite{virtualbox_guestcontrol, virtualbox_guestadditions}.

La etapa de almacenamiento es idempotente porque cada acción está precedida por una guarda que consulta el estado actual antes de modificarlo \parencite{idempotency_concept}. La primera guarda utiliza \texttt{lsblk} para identificar el disco adicional por su tamaño y comprobar si carece de tabla de particiones; solo en ese caso crea una tabla GPT con una partición mediante \texttt{parted} \parencite{linux_storage_management}. La segunda verifica, con \texttt{blkid}, si la partición ya contiene un sistema de archivos; de lo contrario ejecuta \texttt{mkfs.ext4} o \texttt{mkfs.xfs} según la configuración de la \vm{}. La tercera consulta \texttt{findmnt} y el contenido de \texttt{/etc/fstab} y, únicamente si el montaje no está declarado ni activo, obtiene el UUID, crea el punto de montaje con \texttt{mkdir -p}, añade la entrada por UUID para que el montaje persista entre reinicios y ejecuta \texttt{mount -a} \parencite{uuid_fstab, lsb_fhs}.

Finalmente, la etapa de verificación confirma con \texttt{findmnt} y \texttt{df -h} que el almacenamiento quedó montado y registra el resultado: la ejecución termina con código~0 cuando la verificación es correcta y con código~1 cuando alguna guarda o verificación falla, de modo que el estado inconsistente se detecta de inmediato \parencite{bash_best_practices}. Como consecuencia de las guardas, re-ejecutar la solución sobre una \vm{} ya configurada no reformatea el disco ni duplica entradas en \texttt{/etc/fstab} y deja intactos los datos persistidos; la misma secuencia aplicada a \texttt{VM1} y \texttt{VM2} cambia únicamente los parámetros de entrada \parencite{idempotency_concept, ansible_up_running}.

\begin{figure}[H]
    \centering
    \includegraphics[width=0.7\textwidth]{algoritmo.png}
    \caption{Diagrama de procesos de la automatización idempotente, con las etapas ejecutadas en el \host{} y en el sistema invitado}
    \label{fig:algoritmo}
\end{figure}